\documentclass[12pt,a4paper]{article}
\usepackage{amsmath}
\usepackage{amssymb}
\usepackage{amsthm}
\usepackage{mathtools}
\usepackage{geometry}
\usepackage{hyperref}
\usepackage{xcolor}
\usepackage{enumitem}

\geometry{margin=1in}
\emergencystretch=3em

\theoremstyle{definition}
\newtheorem{definition}{Definition}[section]
\newtheorem{theorem}{Theorem}[section]
\newtheorem{lemma}[theorem]{Lemma}
\newtheorem{proposition}[theorem]{Proposition}
\newtheorem{remark}{Remark}[section]

\newcommand{\Nat}{\mathbb{N}}
\newcommand{\Real}{\mathbb{R}}
\newcommand{\Complex}{\mathbb{C}}
\newcommand{\K}{\mathbb{K}}
\newcommand{\norm}[1]{\left\|#1\right\|}
\newcommand{\spanop}{\operatorname{span}}
\newcommand{\closure}{\operatorname{closure}}
\newcommand{\HasSum}{\operatorname{HasSum}}
\newcommand{\Finset}{\operatorname{Finset}}
\newcommand{\range}{\operatorname{range}}
\newcommand{\atTop}{\operatorname{atTop}}

\title{\textbf{Unconditional Schauder Bases and Finite Sign Criteria}\\
  \large Current Lean 4 Formalization}
\author{Daniel Smania}
\date{May 2026}

\begin{document}

\maketitle

\begin{abstract}
This document describes the current Lean 4 formalization in this project.  The
development defines ordered Schauder bases, unconditional Schauder bases indexed
by $\Nat$, an
abstract-index version based on Lean's unconditional finite-set summation
filter, and a finite sign criterion for constructing unconditional Schauder
bases.  The main public existence theorems are
listed in Section~\ref{sec:main-theorems}.
\end{abstract}

\tableofcontents

\newpage

\section{Introduction}

\subsection{Purpose of the Project}

The project formalizes a reusable criterion for producing unconditional
Schauder bases in complete normed spaces.  It is designed as a small Lean
library that downstream projects can use once they have proved three
mathematical facts about a family of vectors:

\begin{enumerate}
  \item the family has dense closed linear span;
  \item every vector in the family is nonzero;
  \item finite signed sums are uniformly bounded by the corresponding unsigned
    finite sums.
\end{enumerate}

The criterion is first stated for an arbitrary index type.  This is the most
natural form in Lean and in applications, because unconditional summability is
controlled by finite subsets and does not require a preferred order on the
index set.  The familiar sequence-indexed theorem is then recovered by taking
the index type to be $\Nat$.

\subsection{Formalized Pipeline}

The current Lean file follows this pipeline:

\begin{enumerate}
  \item Define the ordered partial-sum convergence used by Schauder bases.
  \item Bundle $\Nat$-indexed Schauder bases with continuous coordinate maps.
  \item Define unconditionality using \verb|HasSum|.
  \item Record three rearrangement formulations: \verb|HasSum| reindexing,
    finite-set \verb|Filter.Tendsto| convergence, and the classical criterion
    using ordered partial sums after every permutation of $\Nat$.
  \item Define abstractly indexed unconditional Schauder bases.
  \item Prove that abstract coordinates have the expected Kronecker-delta
    behavior on basis vectors.
  \item Enumerate an abstract basis by an equivalence $\Nat \simeq I$ to obtain
    a usual $\Nat$-indexed unconditional Schauder basis.
  \item Convert a usual unconditional Schauder basis back to the abstract-index
    API with index type $\Nat$.
  \item Convert finite sign bounds into finite projection bounds.
  \item Use the projection bounds to obtain linear independence and continuous
    coordinate maps.
  \item Prove unconditional convergence and uniqueness of the constructed
    expansion.
  \item Package the construction as public abstract-index and $\Nat$-indexed
    existence theorems.
\end{enumerate}

\newpage

\section{Ordered Schauder Bases}

\subsection{Schauder Sums}

\begin{definition}[Ordered Schauder sum]
In Lean, the ordered convergence used by Schauder bases is named
\begin{center}
\verb|HasSchauderSum|.
\end{center}
For a scalar field $\K$, a normed space $E$, a sequence of vectors
$b : \Nat \to E$, coefficients $a : \Nat \to \K$, and a vector $x : E$, it says
\[
  \lim_{N\to\infty}
  \sum_{n\in \Finset.\range(N)} a_n b_n
  =
  x.
\]
Lean expresses this as a \verb|Filter.Tendsto| statement along \verb|atTop|.
\end{definition}

\begin{remark}
This notion is intentionally weaker than \verb|HasSum|.  A Schauder basis is
ordered: the partial sums are taken over initial segments
\verb|Finset.range N|.  By contrast, \verb|HasSum| uses the finite-set filter
and is the appropriate primitive for unconditional summability.
\end{remark}

\subsection{Bundled Schauder Bases}

\begin{definition}[Schauder basis]
The Lean structure
\begin{center}
\verb|SchauderBasis|
\end{center}
contains:
\begin{enumerate}[label=(\roman*)]
  \item basis vectors \verb|basis : Nat -> E|;
  \item continuous coordinate maps \verb|coeff : Nat -> E ->L[K] K|;
  \item a representation theorem saying every vector is the limit of its
    ordered coordinate partial sums;
  \item uniqueness of coefficients for any ordered Schauder expansion.
\end{enumerate}
\end{definition}

The coordinate of a vector $x$ at index $n$ is exposed as
\[
  \verb|SchauderBasis.coord b n x|.
\]

\newpage

\section{Unconditional Schauder Bases}

\subsection{Unconditionality}

\begin{definition}[Unconditional Schauder basis]
For a bundled Schauder basis \verb|b : SchauderBasis K E|, the predicate
\[
  \verb|b.IsUnconditional|
\]
means that every coordinate expansion is a \verb|HasSum|:
\[
  \HasSum\left(n\mapsto b.\operatorname{coeff}_n(x)\,b_n\right)(x)
  \quad\text{for every }x\in E.
\]
\end{definition}

\begin{theorem}[\texorpdfstring{\texttt{HasSum}}{HasSum} reindexing]
The theorem
\begin{center}
\verb|SchauderBasis.isUnconditional_iff_hasSum_rearranged|
\end{center}
states that \verb|HasSum|-based unconditionality is invariant under every
permutation:
\[
  \HasSum\left(
    n\mapsto b.\operatorname{coeff}_{\sigma(n)}(x)\,b_{\sigma(n)}
  \right)(x)
\]
for every permutation $\sigma$ of $\Nat$.
\end{theorem}

\begin{theorem}[Finite-set rearranged convergence]
The theorem
\begin{center}
\verb|SchauderBasis.isUnconditional_iff_tendsto_rearranged|
\end{center}
is the same unconditionality statement with the abbreviation
\verb|HasSum| unfolded.  It says that \verb|b.IsUnconditional| is equivalent to
\[
  \lim_{s\to\infty}
  \sum_{n\in s}
    b.\operatorname{coeff}_{\sigma(n)}(x)b_{\sigma(n)}
  =
  x,
\]
where $s$ ranges over finite subsets of $\Nat$ along the finite-set
\verb|atTop| filter.
\end{theorem}

\begin{theorem}[Ordered rearranged convergence]
The one-way theorem
\begin{center}
\verb|SchauderBasis.isUnconditional_tendsto_rearranged|
\end{center}
states that unconditionality implies convergence of every permuted ordered
expansion.  This is often the most convenient forward-use form of the
classical rearrangement statement.
\end{theorem}

\begin{theorem}[Ordered rearrangement criterion]
The theorem
\begin{center}
\verb|SchauderBasis.isUnconditional_iff_tendsto_ordered_rearranged|
\end{center}
states the classical ordered partial-sum criterion:
\verb|b.IsUnconditional| is equivalent to
\[
  \lim_{N\to\infty}
  \sum_{n\in\Finset.\range(N)}
    b.\operatorname{coeff}_{\sigma(n)}(x)b_{\sigma(n)}
  =
  x
\]
for every vector $x$ and every permutation $\sigma$ of $\Nat$.  This is the
direct formal version of the usual statement that every rearranged Schauder
expansion converges to the same vector.
\end{theorem}

\begin{definition}[Bundled unconditional Schauder basis]
The Lean structure
\begin{center}
\verb|UnconditionalSchauderBasis|
\end{center}
bundles a \verb|SchauderBasis| together with the unconditionality predicate
above.  It also provides convenient projections
\verb|UnconditionalSchauderBasis.basis| and
\verb|UnconditionalSchauderBasis.coeff|.
\end{definition}

\newpage

\section{Abstractly Indexed Bases}

\subsection{Motivation}

Many applications naturally produce a family indexed by a type that is not
literally $\Nat$.  For such a family there is no canonical order of partial
sums.  The right primitive is therefore unconditional convergence over finite
subsets, expressed in Lean by \verb|HasSum|.

\begin{definition}[Abstract-index unconditional Schauder basis]
The structure
\begin{center}
\verb|UnconditionalSchauderBasisAbstractIndex|
\end{center}
is indexed by an arbitrary type $I$.  It contains:
\begin{enumerate}[label=(\roman*)]
  \item basis vectors \verb|basis : I -> E|;
  \item continuous coordinate maps \verb|coeff : I -> E ->L[K] K|;
  \item a \verb|HasSum| representation theorem for every vector;
  \item uniqueness of coefficients for every unconditional expansion.
\end{enumerate}
\end{definition}

\subsection{Coordinates of Basis Vectors}

From uniqueness of unconditional expansions, the file proves the standard
Kronecker-delta behavior.

\begin{lemma}[Self coordinate]
For an abstract-index unconditional Schauder basis $b$,
\[
  b.\operatorname{coeff}_i(b.\operatorname{basis}_i)=1.
\]
The Lean theorem is \verb|coeff_basis_self|.
\end{lemma}

\begin{lemma}[Off-diagonal coordinate]
If $j\neq i$, then
\[
  b.\operatorname{coeff}_i(b.\operatorname{basis}_j)=0.
\]
The Lean theorem is \verb|coeff_basis_ne|.
\end{lemma}

\subsection{Enumeration}

\begin{theorem}[Enumerating an abstract basis]
Given an abstract-index unconditional Schauder basis indexed by $I$ and an
equivalence $e : \Nat \simeq I$, Lean constructs a usual $\Nat$-indexed
unconditional Schauder basis:
\begin{center}
\verb|UnconditionalSchauderBasisAbstractIndex.toUnconditionalSchauderBasis|.
\end{center}
The resulting sequence has basis vectors $n\mapsto b.\operatorname{basis}(e n)$
and coordinate maps $n\mapsto b.\operatorname{coeff}(e n)$.
\end{theorem}

The file also proves that this enumeration does not change the set of basis
vectors or the set of coordinate maps:
\[
\begin{aligned}
  &\verb|range_basis_toUnconditionalSchauderBasis|,\\
  &\verb|range_coeff_toUnconditionalSchauderBasis|.
\end{aligned}
\]

\subsection{Forgetting the Order}

Conversely, every usual unconditional Schauder basis gives an abstractly
indexed one over $\Nat$:
\begin{center}
\verb|UnconditionalSchauderBasis.toUnconditionalSchauderBasisAbstractIndex|.
\end{center}
The existence theorem
\[
  \verb|exists_unconditionalSchauderBasisAbstractIndex_nat|
\]
records that the basis vectors and coordinate maps are unchanged.

\newpage

\section{Finite Sign Criterion}

\subsection{Dense Span}

\begin{definition}[Dense closed span]
For a family $x : I \to E$, the predicate
\begin{center}
\verb|UnconditionalCriterion.HasDenseSpan x|
\end{center}
means
\[
  \closure\bigl(\spanop_{\K}\{x_i : i\in I\}\bigr)=E.
\]
In Lean this is expressed as equality with \verb|Set.univ| after coercing the
span to a set.
\end{definition}

\subsection{Finite Sign Bound}

\begin{definition}[Finite sign estimate]
For a family $x : I\to E$ and a constant $C\in\Real$, the predicate
\begin{center}
\verb|UnconditionalCriterion.HasFiniteSignBound x C|
\end{center}
says that for every finite set $s\subset I$, every coefficient function
$a : I\to\K$, and every sign function $\varepsilon : I\to\K$ with
$\varepsilon_i=1$ or $\varepsilon_i=-1$ on $s$,
\[
  \norm{\sum_{i\in s}(\varepsilon_i a_i)x_i}
  \le
  C\,\norm{\sum_{i\in s}a_i x_i}.
\]
\end{definition}

\begin{remark}
The scalar field assumptions in the criterion are those in the Lean file:
\verb|NontriviallyNormedField K|, \verb|CharZero K|, and
\verb|CompleteSpace K|, together with a complete normed space $E$ over $K$.
The characteristic-zero assumption is used for the signs and the projection
argument involving the scalar $2$.
\end{remark}

\newpage

\section{Proof Strategy}

\subsection{From Signs to Projections}

The first step converts finite sign estimates into uniform bounds for finite
coordinate projections.  Given finite sets $t\subseteq s$, choose signs
\[
  \varepsilon_i =
  \begin{cases}
    1, & i\in t,\\
    -1, & i\notin t.
  \end{cases}
\]
The algebraic identity
\[
  \sum_{i\in s}(\varepsilon_i a_i)x_i
  =
  2\sum_{i\in t}a_i x_i
  -
  \sum_{i\in s}a_i x_i
\]
leads to the projection estimate
\[
  \norm{\sum_{i\in t}a_i x_i}
  \le
  \norm{2^{-1}}\,(C+1)\,
  \norm{\sum_{i\in s}a_i x_i}.
\]
The Lean predicate for this intermediate statement is private and named
\begin{center}
{\small\verb|FiniteProjectionBound|}
\end{center}

\subsection{Linear Independence and Coordinate Maps}

The projection bound implies linear independence of the family $x$, provided
each $x_i$ is nonzero.  This is the lemma
\begin{center}
\verb|linearIndependent_of_finiteProjectionBound|.
\end{center}

Once linear independence is available, the proof defines algebraic coordinate
maps on the span of the family.  The projection bound gives the norm control
needed to extend those maps continuously from the dense span to all of $E$.
This step is packaged in
\begin{center}
\verb|exists_coordMaps_of_finiteProjectionBound|.
\end{center}

\subsection{Convergence}

Let $P_s$ be the finite partial-sum projection associated to a finite set
$s\subset I$:
\[
  P_s(y)=\sum_{i\in s}\operatorname{coeff}_i(y)x_i.
\]
On algebraic finite linear combinations, $P_s$ is eventually exact once $s$
contains the support of the combination.  The uniform projection bound then
extends convergence from the dense algebraic span to every vector in $E$.

Lean proves this analytic core as
\begin{center}
\verb|coordMaps_tendsto_finite_partial_sums_of_finiteProjectionBound|.
\end{center}
It then repackages the statement as a \verb|HasSum| representation theorem.

\subsection{Uniqueness}

Uniqueness of coefficients follows by applying the constructed coordinate
functionals to an arbitrary unconditional expansion.  Once the finite set
contains the chosen index $i$, the $i$th coordinate of the finite partial sum
is exactly the coefficient $a_i$.  Passing to the \verb|HasSum| limit identifies
$a_i$ with the constructed coordinate of the sum.

\newpage

\section{Main Theorems}
\label{sec:main-theorems}

\subsection{Abstract-Index Criterion}

\begin{theorem}[Abstract-index finite sign criterion]
Assume $x : I\to E$ has dense closed linear span, every $x_i$ is nonzero, and
$x$ satisfies a finite sign bound with constant $C\ge 0$.  Then there exists an
abstractly indexed unconditional Schauder basis whose basis family is exactly
$x$.

The Lean theorem is
\begin{center}
{\scriptsize\ttfamily
UnconditionalCriterion.\\
exists\_unconditionalSchauderBasisAbstractIndex\_of\_finiteSignBound}
\end{center}
\end{theorem}

\subsection{\texorpdfstring{$\Nat$}{Nat}-Indexed Criterion}

\begin{theorem}[Sequence-index finite sign criterion]
Assume $x : \Nat\to E$ has dense closed linear span, every $x_n$ is nonzero,
and $x$ satisfies a finite sign bound with constant $C\ge 0$.  Then there
exists a usual unconditional Schauder basis whose basis sequence is exactly
$x$.

The Lean theorem is
\begin{center}
{\scriptsize\ttfamily
UnconditionalCriterion.\\
exists\_unconditionalSchauderBasis\_of\_finiteSignBound}
\end{center}
\end{theorem}

The sequence-index theorem is obtained from the abstract-index theorem by
specializing the index type to $\Nat$ and using the canonical equivalence
\verb|Equiv.refl Nat|.

\newpage

\section{File Guide}

The project currently has one main Lean file:

\begin{itemize}
  \item \verb|UnconditionalSchauderBasis.lean|:
    ordered Schauder sums, bundled Schauder bases, unconditional
    Schauder bases, abstract-index unconditional bases, conversions between
    the two APIs, and the complete finite sign criterion.
\end{itemize}

The repository also contains:

\begin{itemize}
  \item \verb|README.md|:
    short project overview and build instructions;
  \item \verb|docpdf/Documentation.tex| and
    \verb|docpdf/Documentation.pdf|:
    narrative PDF documentation for the current formalization;
  \item \verb|lakefile.toml|:
    Lake package configuration;
  \item \verb|lean-toolchain|:
    Lean toolchain pin;
  \item \verb|lake-manifest.json|:
    resolved dependency manifest.
\end{itemize}

\section{Conclusion}

The current formalization provides a compact, reusable bridge from finite sign
estimates to unconditional Schauder bases.  The abstract-index theorem is the
preferred general API for downstream work, while the $\Nat$-indexed theorem
keeps the familiar sequence-based Schauder-basis interface available when an
ordered basis is wanted.

\end{document}
